% !TeX program = lualatex
\documentclass[a4paper,11pt]{ltjsarticle}
\usepackage{icat-common}
\usepackage{lmodern}
\usepackage{mathtools}
\usepackage{booktabs}
\usepackage{array}
\usepackage[unicode]{hyperref}

\theoremstyle{definition}
\newtheorem{definition}{定義}[section]
\newtheorem{axiom}[definition]{公理}
\newtheorem{remark}[definition]{注記}
\newtheorem{example}[definition]{例}
\theoremstyle{plain}
\newtheorem{proposition}[definition]{命題}
\newtheorem{theorem}[definition]{定理}

\newcommand{\Fzero}{F_{0}}
\newcommand{\CircleRef}{\bigcirc^{\mathrm{ref}}}
\newcommand{\CircleEval}{\bigcirc^{\mathrm{eval}}}
\newcommand{\Obj}{\mathcal M}
\newcommand{\Obs}{\mathcal R}
\newcommand{\Time}{\mathcal T}
\newcommand{\Val}{\mathcal V}
\newcommand{\Set}{\mathbf{Set}}
\newcommand{\Lan}{\operatorname{Lan}}
\newcommand{\Ran}{\operatorname{Ran}}
\newcommand{\Nat}{\operatorname{Nat}}
\newcommand{\op}{\mathrm{op}}
\newcommand{\Chu}{\mathbf{Chu}}
\newcommand{\PSh}{\widehat{\mathcal C}}
\newcommand{\CoPSh}{\check{\mathcal C}}
\newcommand{\QRB}{\mathrm{QRB}}
\newcommand{\Res}{\mathrm{Res}}
\newcommand{\supp}{\mathrm{supp}}
\newcommand{\Spc}{\mathrm{Spc}}
\newcommand{\Ctx}{\mathbf{Ctx}}
\newcommand{\Ty}{\mathrm{Ty}}
\newcommand{\Tm}{\mathrm{Tm}}

\title{意味と指示の構造論\\
{\large ---直交性、鏡像構成、カン拡張、および指示の操作的先行性---}}
\author{千葉将臣}
\date{2026-07-29}

\begin{document}
\maketitle

\begin{abstract}
本稿は、意味と指示の関係を、意味論上の同一性または分離可能性の二択としてではなく、観測対象と観測者像のあいだに成立する構造として分析する。出発点は、通常の圏に現れる対象構造と、カン拡張に現れる文脈変換を、意味と指示の二方向として読むことである。この着想を二つの相補的視点へ分解する。第一の\textbf{直交性の視点}は、意味／観測対象と指示／観測者像を、独立に変化しうる二方向として分析的に区別する。第二の\textbf{鏡像構成の視点}は、両方向が評価関係を介して左右反転し、相互に表現を与える仕組みを捉える。後者の形式的模型としてChu空間を、両側の表現を相互に再構成する形式としてIsbell双対を位置づけ、カン拡張を観測文脈の変更に沿う普遍的な延長・補完として解釈する。

指示方向の標準的な圏論的基盤として、Cartmellのcontextual categoryおよびVoevodskyのC-systemを導入する。文脈と代入の圏を基底とし、各文脈上の型と項をファイバーとして配置すると、指示／観測文脈は基底方向、意味はその上のファイバー方向として形式化される。代入 $f^*$ は体系内観測位置の変更を表し、その左右随伴 $\Sigma_f\dashv f^*\dashv\Pi_f$ および関手圏上の $\Lan_K\dashv K^*\dashv\Ran_K$ は、文脈変更に沿う意味・指示像の再構成を与える。

さらに、単一のChu鏡面を、四重残余束（QRB）に現れる四つの比較操作へ分解する。四重結像とは四つの像が一致することではなく、同一性截断、二重否定閉包、随伴単位、自己参照比較という四方向の結像障害が同じ素点で同時に検出されることである。QRBの外点候補領域は四残余の台の共通部分として与えられ、真理保存性はこの四重構造そのものではなく、その後に確認される意味論的一断面として位置づけられる。

観測者と時間は同一視しない。観測者は区別を行う位置であり、時間は状態および指示点を順序づける構造である。$F_0$ の記号 $\bigcirc$ は、評価値域と現実的指示点という異なる型の役割を担いうるため、本稿では $\CircleEval$ と $\CircleRef$ に型分けする。無点化された鏡像構成そのものは対象側と観測側に優先順位を与えないが、$\CircleRef$ が特定の観測位置を選ぶとき、その位置は観測された意味の成立に操作上先行する。本稿はこれを\textbf{指示先行公理}として定式化する。数学は時間構造と指示点間の関係を記述できるが、無点化された構造だけから現実の「ここ・いま・これ」を選出することはない。これらの主張は数学上の一般的不可能性定理ではなく、構造記述と現実的指示を分離する $F_0$ メタ分析上の原理である。

時間については、文脈の依存順序、項の計算・簡約順序、HoTTのパスが与える変形次元、現実的指示点をもつ観測時間を区別する。これらは相互に関係づけられるが、標準的には同一の時間構造ではない。本稿で「観測者は時間ではない」というとき、主として第四の観測時間と観測文脈を区別している。
\end{abstract}

\tableofcontents

\section{導入}

\subsection{問題設定}

Frege以来、意味と指示の区別は、記号が対象を指すことと、その対象が与えられる仕方を区別するための基本的枠組みである\cite{Frege1892Sinn}。しかし、意味と指示を二つの完成した対象として並置すると、両者を結ぶ作用が第三項として必要になる。他方、両者を同一化すると、同じ対象を異なる仕方で与える構造を失う。

本稿はこの問題を、次の四項の関係として読み替える。
\begin{enumerate}
 \item 意味を担う\textbf{観測対象}、
 \item 指示を成立させる\textbf{観測位置}、
 \item 両者を結ぶ\textbf{評価関係}、
 \item 観測が「ここ・いま・これ」として現実化する\textbf{指示点}。
\end{enumerate}
観測位置は観測者そのものではなく、体系内に表現された観測者像である。体系外の観測者を内部対象として固定すれば、その固定を行う観測作用が再び外部に残るからである。この区別は、$F_0=\{1,0,\infty,\bigcirc\}$ を対象評価層 $\{1,0\}$ と観測者反省層 $\{\infty,\bigcirc\}$ に分ける既存のメタ分析を継承する\cite{ChibaF02026}。

\subsection{本稿の主張}

本稿の中心的主張は次の八点である。
\begin{enumerate}
 \item 意味と指示の\textbf{直交性}は、両者の無関係性ではなく、独立変化可能性を表す。
 \item 意味と指示の\textbf{鏡像構成}は、直交的に区別された二方向が評価関係によって相互規定されることを表す。
 \item contextual category／C-systemは、指示を文脈と代入の基底軸として、意味をその上の型・項のファイバーとして形式化する。
 \item Chu空間は鏡像構成の静的模型を、Isbell双対は両側の表現の相互再構成を与える。
 \item QRBの○は単一の鏡像ではなく、四つの比較操作の残余が同一点で重なる四重結像として捉えられる。
 \item カン拡張は観測者の実在的移動ではなく、体系内の観測文脈の変換に沿う像の延長・補完を与える。
 \item 観測者と時間は異なる。両者は指示点 $\CircleRef$ において交差するが、同一ではない。
 \item 観測された意味が成立するには、$\CircleRef$ による観測位置の指定が操作上先行する。
\end{enumerate}

\subsection{非主張と方法上の注意}

本稿は、$F_0$ とChu空間またはIsbell双対が数学的に同一であるとは主張しない。また、「通常の圏」と「カン拡張の方向の圏」という二種類の圏が、既存の圏論において厳密に直交するとも主張しない。カン拡張は関手圏における普遍構成であり、随伴、極限、余極限と深く関係する\cite{Kan1958,MacLane1998,Riehl2016}。本稿の対応は、$F_0$ の意味論的区別を圏論的に精密化するための構造上の比較である。

\section{意味と指示の二方向}

\subsection{対象側と観測側}

意味構造の領域を $\Obj$、体系内の観測位置の領域を $\Obs$ とする。概念的には、$m\in\Obj$ は観測される対象、$r\in\Obs$ は対象を区別する観測文脈を表す。

\begin{definition}[意味方向]
意味方向とは、対象の構造、対象間の射、同型、双対、随伴など、何がどの構造をもつかを記述する方向である。$F_0$ では主として対象評価層 $\{1,0\}$ に対応する。
\end{definition}

\begin{definition}[指示方向]
指示方向とは、対象がどの文脈から「これ」として与えられるか、および文脈変換後に指示像がどう再構成されるかを記述する方向である。$F_0$ では主として観測者反省層 $\{\infty,\bigcirc\}$ に対応する。
\end{definition}

ここで「指示＝観測者」という略記は、実在的主体との同一化ではない。正確には、指示方向は\textbf{体系内に表現された観測者の位置および指示作用}である。

\subsection{直交性の視点}

\begin{definition}[構造的直交性]
意味方向 $\Obj$ と指示方向 $\Obs$ が構造的に直交するとは、少なくとも概念上、次の二種類の変換を区別できることをいう。
\begin{enumerate}
 \item 対象構造を固定したまま観測文脈を変換すること。
 \item 観測文脈を固定したまま対象構造を変換すること。
\end{enumerate}
\end{definition}

この定義における直交性は、内積が零であるという幾何学的主張ではない。また、両方向のあいだに関係がないことも意味しない。直交性が示すのは、両者を同一の変数へ還元してはならないという分析上の独立性である。

\begin{proposition}[直交性の非分離性]
意味方向と指示方向が構造的に直交していても、両者が完全に分離可能であるとは限らない。
\end{proposition}

\begin{proof}
対象 $m$ がどのように現れるかは観測文脈 $r$ に依存し、観測文脈 $r$ が何を区別する文脈であるかは対象領域 $\Obj$ に依存する。したがって変換方向は区別できるが、具体的な観測結果は両者の組に依存する。
\end{proof}

\subsection{鏡像構成の視点}

直交性が両方向を分解して見る視点であるのに対し、鏡像構成は両方向の結合を見る。評価値域を $\Val$ とし、
\begin{equation}
 \varepsilon:\Obj\times\Obs\longrightarrow\Val
 \label{eq:evaluation}
\end{equation}
を置く。$\varepsilon(m,r)$ は、観測文脈 $r$ において対象 $m$ がどのように現れ、区別されるかを表す。

評価表を転置すれば、
\begin{equation}
 \varepsilon^{\mathsf T}:\Obs\times\Obj\longrightarrow\Val,
 \qquad
 \varepsilon^{\mathsf T}(r,m)=\varepsilon(m,r)
\end{equation}
を得る。これは対象と観測位置が同一になることではなく、同じ評価関係を反対側から読むことである。

\begin{theorem}[二視点の相補性]
直交性と鏡像構成は、意味と指示の異なる現象ではなく、同じ評価構造の分析面と総合面である。
\end{theorem}

\begin{proof}
直交性は式\eqref{eq:evaluation}の二つの入力 $\Obj$ と $\Obs$ を型として区別する。鏡像構成は、それらが評価写像 $\varepsilon$ において結合され、転置によって役割を交換できることを示す。前者は入力の非同一性を、後者は入力間の相互規定性を記述するため、両者は相補的である。
\end{proof}

\section{文脈圏による指示軸の形式化}

\subsection{Contextual categoryとC-system}

依存型理論では、型や項は無文脈に存在するのではなく、判断文脈 $\Gamma$ に相対化される。基本判断は
\begin{equation}
 \Gamma\;\mathsf{ctx},
 \qquad
 \Gamma\vdash A\;\mathsf{type},
 \qquad
 \Gamma\vdash a:A
 \label{eq:contextual-judgments}
\end{equation}
である。Cartmellのcontextual categoryは、文脈、文脈拡張、代入を圏論的に組織するための構造であり、Voevodskyはその構成的な変種をC-systemとして研究した。とりわけ、宇宙圏 $(\mathcal C,p)$ からC-system $CC(\mathcal C,p)$ を構成する手続きが与えられている\cite{Voevodsky2015CSystem}。

Contextual categoryまたはC-systemの対象は文脈 $\Gamma$、射
\begin{equation}
 f:\Delta\longrightarrow\Gamma
\end{equation}
は $\Gamma$ の変数を $\Delta$ 上の項によって置き換える代入として読まれる。型 $A$ による文脈拡張は
\begin{equation}
 \Gamma.A\xrightarrow{\pi_A}\Gamma
\end{equation}
と表され、代入 $f$ に沿う型の再添字付けは、選択された引戻しによって
\begin{equation}
 f^*A\in\Ty(\Delta)
\end{equation}
を与える。こうした文脈上の型理論的構造は、contextual categories、categories with families、split type-categoriesなど複数の同値または近縁な形式で研究されている\cite{AhrensLumsdaineVoevodsky2018}。

\begin{remark}[確立された構造と本稿の解釈]
標準理論で確立されているのは「指示の圏」という名称ではなく、文脈と代入の圏である。本稿は、この文脈圏を指示／観測者像の軸として読む。したがって、以下の対応はcontextual categoryの定義そのものではなく、その $F_0$ 的解釈である。
\end{remark}

\subsection{基底としての指示、ファイバーとしての意味}

文脈の圏を $\Ctx$ とする。各文脈 $\Gamma$ における型の圏または型の集まりを $\Obj(\Gamma)$ と書けば、意味構造は概念的に反変な添字付け
\begin{equation}
 \Obj:\Ctx^{\op}\longrightarrow\mathbf{Cat}
 \label{eq:indexed-meaning}
\end{equation}
として表せる。集合値の簡略化では $\Obj:\Ctx^{\op}\to\Set$ としてもよい。文脈射 $f:\Delta\to\Gamma$ に対して
\begin{equation}
 f^*:\Obj(\Gamma)\longrightarrow\Obj(\Delta)
\end{equation}
が代入・再添字付けを与える。

\begin{definition}[文脈的指示軸]
文脈圏 $\Ctx$ とその射による方向を、体系内で指示位置を組織する\textbf{文脈的指示軸}と呼ぶ。各文脈 $\Gamma$ 上のファイバー $\Obj(\Gamma)$ を、その指示位置において成立する\textbf{意味ファイバー}と呼ぶ。
\end{definition}

このとき、これまでの意味と指示の直交性は、単なる二変数の独立性より精密に、基底方向とファイバー方向の区別として表現される。
\begin{equation}
 \begin{array}{ccl}
 \Gamma\in\Ctx
 &:& \text{指示・観測文脈},\\
 A\in\Obj(\Gamma)
 &:& \text{文脈 $\Gamma$ で成立する意味},\\
 a\in\Tm_{\Gamma}(A)
 &:& \text{その意味の下で指示される項},\\
 f:\Delta\to\Gamma
 &:& \text{体系内観測位置の変換}.
 \end{array}
 \label{eq:context-f0-correspondence}
\end{equation}

\begin{proposition}[基底・ファイバーとしての直交性]
文脈的指示軸と意味方向の直交性は、基底 $\Ctx$ に沿う移動と、各ファイバー $\Obj(\Gamma)$ 内の変換を型の異なる操作として区別することによって形式化される。
\end{proposition}

\begin{proof}
文脈変更 $f:\Delta\to\Gamma$ は異なるファイバー間の再添字付け $f^*$ を誘導する。一方、固定した $\Gamma$ における射 $A\to B$ は同じファイバー $\Obj(\Gamma)$ 内の意味変換である。前者は基底方向、後者は垂直方向であり、同じ射型ではない。ただし再添字付けがファイバー内構造へ作用するため、両方向は独立に区別されながら相互に関係する。
\end{proof}

この構図は、「指示＝観測者」という略記をさらに限定する。$\Ctx$ が表現するのは実在する観測者ではなく、体系内で利用可能な観測文脈の全体である。したがってcontextual categoryは観測者そのものではなく、観測者像を形式化する。

\subsection{代入の随伴とカン拡張}

適切な comprehension category、局所デカルト閉圏、または同等の型理論的意味論では、再添字付け $f^*$ の左右に
\begin{equation}
 \Sigma_f\dashv f^*\dashv\Pi_f
 \label{eq:dependent-adjunctions}
\end{equation}
が存在する。$\Sigma_f$ は依存和に、$\Pi_f$ は依存積に対応し、文脈変更に沿って意味を前向きまたは制約的に再構成する。

関手圏において文脈関手 $K:\mathcal C\to\mathcal D$ を用いる場合には、制限関手 $K^*$ の左右随伴として
\begin{equation}
 \Lan_K\dashv K^*\dashv\Ran_K
 \label{eq:context-kan-triple}
\end{equation}
を得る。式\eqref{eq:dependent-adjunctions}は文脈圏の各射に沿う局所的・依存的な再構成を、式\eqref{eq:context-kan-triple}は関手圏全体における文脈変更を記述する。この区別により、カン拡張を実在的観測者の移動ではなく、文脈圏上のデータを別の基底へ移送する普遍構成として位置づけられる。

\begin{remark}[Kan fibrationとの区別]
単体的集合によるUnivalent Foundationsのモデルで用いられるKan fibrationと、本稿で扱うKan extensionは別の概念である。単体的モデルがKan fibrationを用いることから、指示軸としてのKan extensionが直接導かれるわけではない。カン拡張への接続は、文脈関手に沿う制限 $K^*$ とその左右随伴から得られる。
\end{remark}

Kapulkin--Lumsdaineによる単体的モデルは、宇宙を用いて依存型理論のcontextual categoryを構成し、単体的集合上のunivalent universeを与える。その帰結として、contextual categoryの言葉で定式化された単一のunivalent universeをもつMartin-Löf型理論が、二つの到達不能基数をもつZFCと少なくとも同程度に無矛盾であることが示される\cite{KapulkinLumsdaine2021}。これは文脈圏がHoTTの一つの標準的な意味論的核であることを示すが、文脈圏そのものを時間や実在的観測者と同一化するものではない。

\subsection{文脈圏上の四重結像}

文脈圏を導入すると、QRBの四比較操作も文脈相対化できる。各 $\Gamma\in\Ctx$ の意味ファイバーに
\begin{equation}
 M_i^{\Gamma}:\Obj(\Gamma)\longrightarrow\Obj(\Gamma),
 \qquad
 i\in\{S,N,SN,NS\}
\end{equation}
を置く。文脈変更 $f:\Delta\to\Gamma$ に対して、四重結像の安定性は比較写像
\begin{equation}
 \beta_{i,f}:
 f^*M_i^{\Gamma}
 \Longrightarrow
 M_i^{\Delta}f^*
 \label{eq:contextual-fourfold-comparison}
\end{equation}
によって測られる。

\begin{definition}[文脈安定な四重結像]
すべての文脈射 $f$ と四操作 $i$ について、式\eqref{eq:contextual-fourfold-comparison}が定義され同型となり、恒等射および合成と整合するとき、四重結像は文脈安定であるという。
\end{definition}

これは、同一性截断、二重否定閉包、随伴単位、自己参照比較を、観測文脈を変更する前に適用しても後に適用しても同じ結果になるというBeck--Chevalley型の条件である。比較が非可逆となる場合、その残余は単なる対象内部の障害ではなく、文脈変更と意味操作の不整合を記録する。

適切な安定ファイバーと台の理論が与えられれば、各文脈 $\Gamma$ における外点候補領域
\begin{equation}
 \mathcal O_{X,\Gamma}
 =
 \bigcap_i
 \supp\bigl(\Res_i^{\Gamma}(X)\bigr)
\end{equation}
を考えられる。この構成により、文脈圏はQRBの四重結像がどの観測位置で保存され、どこで同時に崩壊するかを添字付けする基底となる。

\begin{remark}[UniMathとの関係]
UniMathは依存型理論とunivalenceに基づく形式化ライブラリであり、その全体をC-systemの意味論上に構築していると述べるのは強すぎる。一方、contextual categories、C-systems、categories with familiesなどの比較はUniMath上で形式化されており、Voevodskyの型理論メタ理論においてC-systemが主要な道具であったことは明確である\cite{AhrensLumsdaineVoevodsky2018}。
\end{remark}

\section{Chu空間：鏡像構成の静的模型}

\subsection{定義と双対}

固定された集合 $K$ 上のChu空間は、三項組
\begin{equation}
 (A,r,X),
 \qquad
 r:A\times X\longrightarrow K
\end{equation}
として与えられる\cite{Pratt1995}。$A$ を点・状態の側、$X$ を属性・試験の側と読めば、$r$ は状態が各試験に対して示す評価表である。

Chu双対は
\begin{equation}
 (A,r,X)^{\perp}=(X,r^{\mathsf T},A)
\end{equation}
によって両側を交換する。この形は、
\begin{equation}
 (\Obj,\varepsilon,\Obs)^{\perp}
 =
 (\Obs,\varepsilon^{\mathsf T},\Obj)
\end{equation}
という本稿の鏡像構成に対応する。

\begin{remark}[鏡像は同一化ではない]
Chu双対が交換するのは評価関係の左右の役割である。$A=X$ や、点と属性が同一種類の対象であることは要求されない。したがって観測対象と観測者像の交換も、主体と客体の存在論的同一化ではない。
\end{remark}

\subsection{○の型分け}

元の考察では、評価関係を $\varepsilon:\Obj\times\Obs\to\bigcirc$ と略記した。しかし論文上は、評価の値域と指示の発生点を同じ型へ入れるべきではない。そこで本稿では、同じ $F_0$ 記号の型付きの現れとして、
\begin{align}
 \CircleEval &:\text{評価関係が値を取る境界または値域},\\
 \CircleRef &:\text{「ここ・いま・これ」を指定する指示点}
\end{align}
を区別する。Chu空間の $K$ に対応するのは主として $\CircleEval$ であり、後述する点付き観測に対応するのは $\CircleRef$ である。

この型分けは $\bigcirc$ を二つの無関係な記号へ分裂させるものではない。同じメタ記号が、評価面と指示点という異なる位置で現れることを明示し、水準崩壊を防ぐものである。

\section{QRBと四重結像}

\subsection{単一鏡面から四つの比較操作へ}

Chu空間の評価写像は、意味と指示を左右に配置する鏡面の基本形を与える。しかしQRBにおける○は、この単一の鏡面だけから導かれるのではない。QRBでは、体系の意味論的状態を表す対象 $X$ に対して、四つの自己関手
\begin{equation}
 M_i:\mathcal C\longrightarrow\mathcal C,
 \qquad
 i\in\{S,N,SN,NS\}
\end{equation}
と自然変換
\begin{equation}
 \eta_i:\operatorname{Id}_{\mathcal C}\Longrightarrow M_i
\end{equation}
を置く\cite{ChibaQRB2026}。四操作の役割は次の通りである。

\begin{center}
\begin{tabular}{p{0.15\textwidth} p{0.31\textwidth} p{0.43\textwidth}}
\toprule
操作 & QRBでの比較 & 意味‐指示論での結像方向 \\
\midrule
$M_S$ & 同一性の外延的等号への截断 & 構造的自己同一性の結像 \\
$M_N$ & 二重否定層化 & 真理・意味論的閉包の結像 \\
$M_{SN}$ & 随伴から生じるモナド $UF$ & 意味から指示への単位的輸送 \\
$M_{NS}$ & 自己参照比較操作 & 指示から意味への反射・対角化 \\
\bottomrule
\end{tabular}
\end{center}

したがって、本稿の単一評価面 $\CircleEval$ は、QRBとの接続において四つの結像様式へ分解される。ここで「結像」とは $M_iX$ という像だけを指すのではなく、元の状態 $X$ とその像 $M_iX$ を比較する写像 $\eta_i(X)$ を含む。

\subsection{四残余と四重結像}

各比較写像の残余を
\begin{equation}
 \Res_i(X)
 =
 \operatorname{cofib}\bigl(\eta_i(X):X\to M_iX\bigr)
\end{equation}
とする。残余が零であることは、適切な安定圏の仮定の下で比較写像が同型であることに対応し、非零残余は元の状態と結像との差異を記録する。

\begin{definition}[四重結像]
対象 $X$ の四重結像とは、四つの比較族
\begin{equation}
 \bigl\{\eta_S(X),\eta_N(X),
 \eta_{SN}(X),\eta_{NS}(X)\bigr\}
\end{equation}
を同時に考察することである。四重結像は四つの像が互いに等しいことを意味せず、四方向の比較が同じ対象 $X$ において重ねて評価されることを意味する。
\end{definition}

四つの残余を束ねたQRBは
\begin{equation}
 \QRB(X)
 =
 \Res_S(X)\oplus\Res_N(X)
 \oplus\Res_{SN}(X)\oplus\Res_{NS}(X)
\end{equation}
である。その外点候補領域は、Balmerスペクトル上で
\begin{equation}
 \mathcal O_X
 =
 \bigcap_{i\in\{S,N,SN,NS\}}
 \supp\bigl(\Res_i(X)\bigr)
 \subseteq
 \Spc(\mathcal C^c)
 \label{eq:qrb-external-locus}
\end{equation}
と定義される。この共通部分に属する点 $\mathfrak p$ では、四つの局所化比較写像 $\eta_i(X)_{\mathfrak p}$ が同時に非同型となる。

\begin{proposition}[四重不動点とQRB外点の対極性]
四方向すべてで
\begin{equation}
 \eta_i(X)\text{ が同型}
 \qquad(i\in\{S,N,SN,NS\})
 \label{eq:fourfold-fixed-point}
\end{equation}
となる状態を四重不動点と呼ぶならば、QRB外点候補は、局所的に四比較写像がすべて非同型となる対極的な障害領域である。
\end{proposition}

したがって、四重結像には二つの極がある。四比較がすべて可逆であれば、同一性、意味論的閉包、随伴輸送、自己参照の全方向で自己同一性が保存される。反対に、それらが同じ点ですべて非可逆となれば、四つの単位的構造が同時に閉じないQRBの外点候補が現れる。

\begin{remark}[候補領域と指示点]
式\eqref{eq:qrb-external-locus}が与える $\mathcal O_X$ は○そのものではなく候補領域である。QRBの指示点 $\CircleRef$ を得るには、$\mathcal O_X\neq\varnothing$ を示し、その内部の点 $\mathfrak p$ を指定する追加データが必要である。強制法の外点と対応させる場合には、さらに比較写像と対応点の構成が必要になる。
\end{remark}

\subsection{真理保存性の位置}

真理保存性は四重結像の全体ではない。とりわけ $M_N$ の意味論的閉包、および $M_S$ の同一性保存と関係する一断面として読めるが、QRBの形式上、それ自体が四残余の一つとして定義されているわけではない。真理保存性は、四つの比較構造と外点指示が設定された後に、特定の意味論またはモデル拡大について確認される帰結である。

したがって、体系の意味‐指示的自己同一性を真理保存性だけから判定することはできない。完全な自己同一性には式\eqref{eq:fourfold-fixed-point}の四条件が必要であり、QRBの○は逆に、四条件の局所的な同時崩壊によって指示される。

\section{カン拡張：観測文脈に沿う再構成}

\subsection{左・右カン拡張}

関手
\begin{equation}
 K:\mathcal C\to\mathcal D,
 \qquad
 F:\mathcal C\to\mathcal E
\end{equation}
を考える。$K$ を文脈の埋め込みまたは変換、$F$ を元の文脈で与えられた意味・指示データと読む。

左カン拡張 $\Lan_KF:\mathcal D\to\mathcal E$ は、自然変換
\begin{equation}
 \eta:F\Longrightarrow (\Lan_KF)K
\end{equation}
を備え、この形の延長のなかで普遍的である。右カン拡張 $\Ran_KF$ は、
\begin{equation}
 \epsilon:(\Ran_KF)K\Longrightarrow F
\end{equation}
を備え、双対な普遍性を満たす\cite{MacLane1998,Riehl2016}。

$F_0$ 的には、左カン拡張を与えられた局所データから像を生成的に延長する側、右カン拡張を全体との整合条件から像を制約的に補完する側として読むことができる。ただし「前進」「後退」は時間方向を意味せず、普遍性の向きを表す比喩である。

\subsection{観測者の移動という比喩}

カン拡張を「観測者の移動に伴う指示の再構成」と呼ぶとき、移動するのは観測者そのものではない。体系が内部に持つ文脈の表現が $K$ によって変換されるのである。

\begin{definition}[体系内観測位置変換]
体系内観測位置変換とは、実在的観測者の運動ではなく、観測文脈を表す圏または添字のあいだの関手 $K:\mathcal C\to\mathcal D$ である。
\end{definition}

この定義により、観測者自身は自らの指示原点にとどまりながら、体系の側では異なる観測者像が射によって関係づけられる。

\subsection{絶対カン拡張}

カン拡張が絶対的であるとは、そのカン拡張が終域 $\mathcal E$ から出る任意の関手によって保存されることをいう。したがって絶対性は「任意の関手に対してカン拡張が存在する」という意味ではない。

$F_0$ 的には、絶対カン拡張を、後続の表現変更を経ても失われない指示再構成として解釈できる。ただし、同型と絶対カン拡張の対応は厳密な同値ではなく、対象側の安定性と文脈変換側の安定性を対照する類比である。

\subsection{カン拡張内部に現れる二方向の対応}

カン拡張は、通常の圏論的構造とは別に置かれた第二の圏ではない。一つのカン拡張
\begin{equation}
 K:\mathcal C\to\mathcal D,
 \qquad
 F:\mathcal C\to\mathcal E
\end{equation}
の内部には、少なくとも二つの役割を分析的に区別できる。第一は、$\mathcal C,\mathcal D,\mathcal E$ の対象・射・同型・双対・随伴などを扱う\textbf{対象構造側}である。第二は、文脈関手 $K$ に沿って $F$ の像を延長し、普遍性によって新しい指示像を定める\textbf{文脈再構成側}である。

この二方向は別々の数学的体系ではなく、同じカン拡張を異なる役割から見たものである。対象構造側は「何がどの構造をもつか」を記述し、文脈再構成側は「その構造が別の文脈でどのように与え直されるか」を記述する。したがって本稿でいう直交性は、カン拡張の外部に二つの圏を並べることではなく、その内部にある\textbf{構造の変換}と\textbf{文脈に沿う再構成}を混同せずに分離する分析的直交性である。

この観点から、対象構造側に現れる概念と、文脈再構成側に現れる概念との対応を次のように整理する。
\begin{center}
\begin{tabular}{p{0.26\textwidth} p{0.29\textwidth} p{0.34\textwidth}}
\toprule
対象構造側 & 文脈・指示側 & 本稿での位置づけ \\
\midrule
同型 & 絶対カン拡張 & 安定性の類比であり同値ではない \\
双対 & 反対圏 & 双対性を実現する操作上の関係 \\
随伴 & 左カン拡張 & 随伴は適切なカン拡張として表現可能 \\
余随伴的方向 & 右カン拡張 & 双対な普遍再構成 \\
\bottomrule
\end{tabular}
\end{center}
この表の各行は、左右の概念が数学的に同値であることを主張しない。同型と絶対カン拡張は異なる水準の安定性を、双対と反対圏は性質とそれを実現する方向反転を、随伴と左右のカン拡張は関係構造と普遍的再構成を、それぞれ対照している。したがってこの表は、カン拡張内部に共存する対象構造方向と文脈再構成方向を $F_0$ 的に可視化した対応表である。

\subsection{四重結像の文脈横断的保存}

四重結像を導入すると、カン拡張の問題は、一つの像を文脈変換後に延長することから、四つの結像様式がそれぞれ文脈変換と両立するかを調べる問題へ精密化される。各 $i\in\{S,N,SN,NS\}$ について、左カン拡張では
\begin{equation}
 \Lan_K(M_iF)
 \qquad\text{と}\qquad
 M_i(\Lan_KF),
 \label{eq:lan-fourfold-comparison}
\end{equation}
右カン拡張では
\begin{equation}
 \Ran_K(M_iF)
 \qquad\text{と}\qquad
 M_i(\Ran_KF)
 \label{eq:ran-fourfold-comparison}
\end{equation}
を比較する。

式\eqref{eq:lan-fourfold-comparison}および式\eqref{eq:ran-fourfold-comparison}のあいだに自然な比較写像が存在することは一般には自動的でない。$M_i$ の保存性、分配則、またはカン拡張との整合条件が追加されなければならない。比較写像が与えられ、それが同型である場合、$i$ 番目の結像は文脈変換を先に行っても後に行っても同じ像を与える。この性質を、$i$ 番目の結像の\textbf{文脈横断的保存}と呼ぶ。

\begin{definition}[四重文脈保存]
四つすべての $i\in\{S,N,SN,NS\}$ について、$M_i$ と選択した左または右カン拡張との比較写像が存在し、かつ同型であるとき、そのカン拡張は四重結像を文脈横断的に保存するという。
\end{definition}

絶対カン拡張は任意の後続関手によって保存されるため、各 $M_i$ がカン拡張の終域上の自己関手として作用する場合、絶対性は四つの結像を後から適用してもカン拡張の普遍性が失われないための強い十分条件となる。これにより、同型と絶対カン拡張の対応は単なる抽象的類比にとどまらず、四つの自己同一性検査を文脈変更後にも保持する安定性として読み直せる。

反対に、四つの比較の一部だけが非可逆なら、その文脈変換は部分的な結像障害を生む。四つすべての比較が同じ局所点で非可逆となる場合には、適切な安定圏、余ファイバー、台の理論を追加することによって、式\eqref{eq:qrb-external-locus}に類するQRB型の共通障害領域を検討できる。ただし、この構成はカン拡張だけから自動的に得られる定理ではなく、四自己関手と比較写像を明示した上で構成すべき課題である。

\section{Isbell双対：両側の表現の相互再構成}

\subsection{前層と余前層}

小圏 $\mathcal C$ に対して、
\begin{equation}
 \PSh=[\mathcal C^{\op},\Set],
 \qquad
 \CoPSh=[\mathcal C,\Set]
\end{equation}
と置く。前層 $P\in\PSh$ は各対象が外部文脈からどのように試されるかを記録する観測プロフィールとして、余前層 $Q\in\CoPSh$ は各対象からどのような作用が出るかを記録する作用プロフィールとして読める。

Yoneda埋め込みを $y(c)=\mathcal C(-,c)$、余Yoneda埋め込みを $y^{\dagger}(c)=\mathcal C(c,-)$ とする。Isbell共役は概略、
\begin{align}
 \mathcal O(P)(c)&=\Nat(P,y(c)),\\
 \mathcal S(Q)(c)&=\Nat(Q,y^{\dagger}(c))
\end{align}
によって、一方のプロフィールから他方のプロフィールを構成する\cite{Isbell1960}。型を整えると、これは
\begin{equation}
 \mathcal O:\PSh
 \rightleftarrows
 \CoPSh^{\op}:\mathcal S
\end{equation}
という随伴として表される。

\subsection{二重共役と反射的閉包}

共役を二度適用すると、一般に標準写像
\begin{equation}
 P\longrightarrow P^{**}
\end{equation}
が得られるが、これは常に同型とは限らない。$P\cong P^{**}$ となる対象は、両側からの表現が整合して反射的に閉じている。

$F_0$ 的には、第一の共役は対象から観測者像を、第二の共役は観測者像から再構成された対象像を与える。二重共役で元に戻らない差は、鏡像化によって追加された整合条件または失われた識別を表す。

\subsection{Chu空間・Isbell双対・カン拡張}

三者の役割は次のように分けられる。
\begin{center}
\begin{tabular}{p{0.22\textwidth} p{0.66\textwidth}}
\toprule
形式 & 本稿での役割 \\
\midrule
Chu空間 & 対象と観測文脈を評価写像の左右に置く鏡の静的な形 \\
Isbell双対 & 左右の表現を共役させ、相互に再構成する反射運動 \\
カン拡張 & 文脈関手に沿ってデータを普遍的に延長・補完する操作 \\
\bottomrule
\end{tabular}
\end{center}
プロ関手を媒介にすればIsbell共役とカン拡張は統一的に扱えるが、本稿ではこの一般化を構成せず、役割の区別にとどめる。

\section{型理論における同一性と時間解釈}

\subsection{点、パス、高次パス}

Homotopy Type Theoryでは、型 $A$ を空間、項 $a,b:A$ を点、同一性型 $a=_A b$ の項をパスとして読む。二つのパス $p,q:a=_A b$ の同一性 $p=q$ は高次パスである\cite{HoTT2013}。

ここで、HoTTにおける集合は任意の型ではない。集合、すなわち $0$-型では、任意の $a,b:A$ に対して同一性型 $a=b$ が命題であり、二つのパス $p,q:a=b$ は同一視される。したがって、同一化の仕方の差を保持するには一般の高次型が必要である。

ユニヴァレンスは、型間の同値を宇宙内の同一性と結びつける。この構造は、「同じ構造をもつが外延的には別に提示された意味体系」を、宇宙内のパスとして扱う手がかりを与える。

\subsection{時間としてのパスという比喩}

パスを時間的変化の完成した履歴として読むことは可能である。しかし、パス変数は物理的時間そのものではなく、変形を表す一般的な次元である。したがって型理論を「時間を静的に表したもの」と呼ぶ場合、正確には次の意味に限定される。
\begin{quote}
型理論は、変化の履歴をパスとして一つの数学的対象へ内在化し、その履歴間の同一性を高次パスとして比較できる。
\end{quote}

意味体系そのものが変化する場合、時間または変化の添字を $t:\Time$ として
\begin{equation}
 A:\Time\longrightarrow\mathcal U
\end{equation}
を置ける。各 $A(t)$ は時点 $t$ における意味体系であり、
\begin{equation}
 a:\prod_{t:\Time}A(t)
\end{equation}
は意味変化に沿う指示の選択または切断である。

\begin{remark}[瞬間と履歴]
特定の $t$ における $p(t)=q(t)$ は瞬間的な比較である。これに対し $p=q$ は履歴全体の整合的な同一性である。「時間を止めたとき同じ」と「二つの時間履歴が同じ」は区別される。
\end{remark}

\subsection{区別すべき四種類の時間}

型理論、文脈圏、計算意味論、および観測論では、順序や変化を表す複数の構造が用いられる。本稿では、これらを一つの時間概念へ同一化せず、次の四種類に分ける。

\begin{center}
\begin{tabular}{p{0.19\textwidth} p{0.31\textwidth} p{0.39\textwidth}}
\toprule
種類 & 形式的な典型 & 本稿での役割 \\
\midrule
依存時間 & 文脈拡張 $\Gamma\to\Gamma.A$、文脈長 & 型・項が成立する論理的依存順序 \\
計算時間 & 簡約列 $e_0\to e_1\to e_2\to\cdots$ & 項の評価・正規化・実行の順序 \\
パス時間 & $p:a=_A b$、高次パス $p=q$ & 状態間の変形と同一化の仕方 \\
観測時間 & $t_0\to t_1\to t_2$ と現在点 $t_0$ & 「ここ・いま」に相対化された観測事象の順序 \\
\bottomrule
\end{tabular}
\end{center}

\begin{definition}[依存時間]
文脈
\begin{equation}
 \Gamma_0
 \longleftarrow
 \Gamma_1
 \longleftarrow
 \Gamma_2
 \longleftarrow\cdots
\end{equation}
において、文脈拡張によって変数・型・仮定が順次依存する順序を\textbf{依存時間}と呼ぶ。これは物理的な前後ではなく、後の判断が前の文脈を必要とするという形成上の順序である。
\end{definition}

Contextual categoryにおける文脈長、基底文脈への射、文脈拡張は依存時間を静的な圏構造として保持する。代入 $f:\Delta\to\Gamma$ は、この依存順序に沿って文脈を読み替える操作である。

\begin{definition}[計算時間]
項または機械状態の簡約・遷移
\begin{equation}
 e_0\longrightarrow e_1\longrightarrow e_2\longrightarrow\cdots
\end{equation}
が与える操作的順序を\textbf{計算時間}と呼ぶ。計算時間は、項がどの型をもつかではなく、その項がどの評価過程を通るかを記述する。
\end{definition}

同じ型をもつ二つの項が異なる計算履歴をもつことも、異なる文脈表現が同じ正規形へ到達することもありうる。したがって依存時間と計算時間には一般に標準的な同一視はない。

\begin{definition}[パス時間]
同一性型の項
\begin{equation}
 p:a=_A b
\end{equation}
を状態 $a$ から状態 $b$ への変形として読むとき、その変形次元を\textbf{パス時間}と呼ぶ。さらに $p=q$ は、二つの変形履歴のあいだの高次変形を表す。
\end{definition}

パス時間は抽象的な変形パラメータであり、計算時間の一ステップや物理的時間の経過を自動的には表さない。計算履歴からパスを構成する場合や、時間発展を型のパスとしてモデル化する場合には、両者を結ぶ追加の解釈写像が必要である。

\begin{definition}[観測時間]
観測事象を順序づける対象 $\Time$ と、その中の現在点
\begin{equation}
 t_0:1\longrightarrow\Time
\end{equation}
からなる点付き時間構造を\textbf{観測時間}と呼ぶ。$F_0$ 的には、現在の観測位置と時点の選択が $\CircleRef$ によって指示される。
\end{definition}

観測時間は、特定の指示点を基準に「現在」「以前」「以後」を組織するための構造である。文脈圏は可能な観測文脈を組織するが、どの文脈と時点が現実の「ここ・いま」であるかは、点付き観測
\begin{equation}
 \sigma:1\longrightarrow\Obs\times\Time
\end{equation}
によって別に指定される。

\begin{proposition}[四時間の非同一性]
依存時間、計算時間、パス時間、観測時間のあいだには、一般に標準的な同一視は存在しない。これらを関係づけるには、簡約からパスへの写像、時間添字付き型族、文脈を時間へ割り当てる関手などの追加構造が必要である。
\end{proposition}

\begin{proof}
依存時間は判断の形成依存、計算時間は操作的遷移、パス時間は同一性または変形、観測時間は現在点をもつ事象順序をそれぞれ記述する。各構造の対象、射、保存すべき関係が異なるため、定義だけから相互の同一性は得られない。追加のモデルがそれらを結ぶ場合に限り、比較写像または関手として関係づけられる。
\end{proof}

$F_0$ における対応は次のようにまとめられる。
\begin{equation}
 \begin{array}{ccl}
 \text{依存時間} &:& \text{文脈的指示軸の形成順序},\\
 \text{計算時間} &:& \text{体系内部の操作的遷移},\\
 \text{パス時間} &:& \text{意味の同一化と高次変形},\\
 \text{観測時間} &:& \CircleRef\text{ に相対化された指示点列}.
 \end{array}
 \label{eq:four-times-f0}
\end{equation}

したがって、「観測者は時間ではない」という主張は、観測文脈 $\Gamma$ と観測時間 $t$ が異なる添字であることを意味する。観測された意味は、必要に応じて
\begin{equation}
 \Obj:\Ctx^{\op}\times\Time\longrightarrow\mathbf{Cat}
 \label{eq:context-time-indexed-meaning}
\end{equation}
という二重添字付けによって表現される。

\section{観測者と時間の分離}

\subsection{三変数の観測構造}

観測者の体系内表現を $\Obs$、時間構造を $\Time$、観測対象を $\Obj$、観測値を $\Val$ とすれば、観測は
\begin{equation}
\operatorname{ev}:\Obs\times\Time\times\Obj
 \longrightarrow\Val
 \label{eq:temporal-evaluation}
\end{equation}
と表せる。

文脈圏による形式化では、$\Obs$ は $\Ctx$ の対象、または観測に用いる文脈の部分圏として具体化される。したがって観測位置 $r\in\Obs$ は文脈 $\Gamma$ に、観測位置の変換は代入射 $f:\Delta\to\Gamma$ に対応する。一方、時間 $\Time$ は文脈圏とは別の添字構造であり、文脈の長さや文脈拡張をそのまま物理的時間と読むことはしない。

\begin{definition}[観測者と時間]
観測者とは、何をどこから区別するかを指定する位置または規則である。時間とは、状態および観測事象を順序づけ、接続する構造である。
\end{definition}

したがって観測者は時間ではない。観測者像を固定したまま時間を変化させることも、同じ時点で観測位置を変化させることも可能である。両者の独立性は、意味と指示の直交性とは別の第三方向を導入する。

\subsection{「ここ・いま」の点付き構造}

式\eqref{eq:temporal-evaluation}だけでは、どの観測位置と時点が現実の「ここ・いま」であるかは指定されない。そこで
\begin{equation}
 \sigma:1\longrightarrow\Obs\times\Time,
 \qquad
 \sigma(*)=(r_0,t_0)
 \label{eq:pointing}
\end{equation}
という点を選ぶ。$F_0$ 的には、この選択作用を $\CircleRef$ として指示する。

\begin{proposition}[指示点の非導出性]
無点化された構造 $(\Obs,\Time,\operatorname{ev})$ は、それだけでは特権的な点 $(r_0,t_0)$ を指定しない。
\end{proposition}

\begin{proof}
特権点は式\eqref{eq:pointing}という追加データである。自己同型によって点が交換されうる構造では、構造不変な仕方で一つの点を選べない場合がある。一般にも、点の選択は無点化構造の定義には含まれていない。
\end{proof}

この命題は、数学が現在時刻を表せないという主張ではない。点 $t_0:1\to\Time$ を追加すれば現在点は表現できる。区別されるのは、時間構造を記述することと、ある点を現実の現在として\textbf{指示すること}である。

\subsection{指示の操作的先行性}

評価写像
\begin{equation}
 \varepsilon:\Obj\times\Obs\longrightarrow\Val
\end{equation}
だけを見れば、$\Obj$ と $\Obs$ はともに入力であり、一方が他方に先行するとはいえない。Chu双対も評価表の左右を交換するため、無点化された鏡像構成は基本的に対称である。指示の先行性は、この対称な構造だけから導かれるのではなく、$\CircleRef$ が一方を現実の観測位置として選ぶことによって生じる。

観測された意味として対象 $m$ が現れることを $\mathsf{Appear}(m)$ と書く。また、$\operatorname{ev}(r,t,m)$ が定義されることを $\operatorname{ev}(r,t,m)\mathord\downarrow$ と書く。

\begin{axiom}[指示先行公理]
対象が観測された意味として現れるならば、その対象を指示する観測位置と時間位置が指定されている。すなわち、
\begin{equation}
 \mathsf{Appear}(m)
 \Longrightarrow
 \exists r:\Obs\;\exists t:\Time\;
 \operatorname{ev}(r,t,m)\mathord\downarrow .
 \label{ax:reference-priority}
\end{equation}
\end{axiom}

式\eqref{ax:reference-priority}の「先行」は時間的順序ではない。意味の現実的な現れを定義するために、観測位置が先に型付けされていなければならないという\textbf{操作的依存順序}である。これを
\begin{equation}
 (r,t):\Obs\times\Time
 \prec_{\mathrm{op}}
 \mathsf{Appear}(m)
\end{equation}
と書く。

この依存関係は、意味の型を観測位置と時間位置に依存する型族
\begin{equation}
 \mathcal M:\Obs\times\Time\longrightarrow\mathcal U
\end{equation}
として精密化すれば、判断
\begin{equation}
 r:\Obs,\;t:\Time
 \;\vdash\;
 m:\mathcal M(r,t)
 \label{eq:dependent-meaning}
\end{equation}
に現れる。式\eqref{eq:dependent-meaning}では、$m$ を観測された意味として型付けする前に、文脈内に $r$ と $t$ が存在しなければならない。

Contextual categoryの判断順序では、この先行性はさらに基本的な形で現れる。判断 $\Gamma\vdash A\;\mathsf{type}$ や $\Gamma\vdash a:A$ を形成するには、まず $\Gamma\;\mathsf{ctx}$ が成立していなければならない。したがって指示先行公理は、文脈形成が型・項判断に先行するという依存型理論の構文的依存関係を、$F_0$ の観測論へ読み替えたものでもある。

\begin{proposition}[指示先行性の弱い解釈]
指示先行公理から導かれるのは、観測された意味の成立が指示位置に依存するという主張であって、観測されない対象または意味が存在しないという主張ではない。
\end{proposition}

\begin{proof}
公理\eqref{ax:reference-priority}の前件は $\mathsf{Appear}(m)$ であり、対象 $m$ の存在そのものではない。したがって同公理は、現れから指示位置の存在を導くが、指示されていない対象の非存在を導かない。
\end{proof}

\begin{remark}[強い存在論的先行性]
「指示されない意味は存在しない」という主張は、
\begin{equation}
 m\text{ が存在する}
 \Longrightarrow
 \exists r,t\;\operatorname{ev}(r,t,m)\mathord\downarrow
\end{equation}
という、より強い存在論的公理を必要とする。本稿はこの公理を採用しない。採用するのは、現実化された観測意味に限定された操作的先行性である。
\end{remark}

$F_0$ 的には、$\CircleEval$ は対象と観測者像を左右対称な評価関係へ置き、$\CircleRef$ はその対称性を破って、特定の観測位置を現実の指示原点として選ぶ。したがって指示の先行性は鏡面そのものからではなく、鏡面の一方を「ここ」として極性化する $\CircleRef$ の作用から生じる。

\subsection{時間性と指示点列}

単独の $\CircleRef$ は時間全体ではなく、時間の切断面を示す。時間性は
\begin{equation}
 \CircleRef_{t_0}\longrightarrow
 \CircleRef_{t_1}\longrightarrow
 \CircleRef_{t_2}\longrightarrow\cdots
\end{equation}
という指示点間の関係として表現される。体系はこの関係を順序、射、パスまたは流れとして記述できる。他方、どの点が現実の現在として指示されているかは、点付きデータに属する。

\section{F0における統合構図}

\subsection{三方向と二つの○}

以上をまとめると、本稿の構造は次の三方向をもつ。
\begin{enumerate}
 \item $\Obj$：意味／観測対象の方向、
 \item $\Obs$：指示／観測者像の方向、
 \item $\Time$：観測事象を接続する時間の方向。
\end{enumerate}
また $\bigcirc$ は型付きで二つの役割をもつ。
\begin{enumerate}
 \item $\CircleEval$：$\Obj$ と $\Obs$ の評価関係が成立する鏡面、
 \item $\CircleRef$：$\Obs$ と $\Time$ の特定の組を「ここ・いま」として選ぶ指示点。
\end{enumerate}

$\CircleEval$ が与える鏡像関係は左右対称であるが、QRBとの接続では、この鏡面は
\begin{equation}
 M_S,\quad M_N,\quad M_{SN},\quad M_{NS}
\end{equation}
の四つの比較方向へ分解される。四重結像の共通残余領域は○の候補領域を与え、$\CircleRef$ はその内部の特定点と観測位置を選択して対称性を極性化する。この選択によって、指示は観測された意味に対して操作的に先行する。

概念図は
\begin{equation}
 \text{意味／観測対象}
 \quad\xleftrightarrow{\ \CircleEval\ }\quad
 \text{指示／観測者像}
\end{equation}
と
\begin{equation}
 \CircleRef_{t_0}\to\CircleRef_{t_1}\to\CircleRef_{t_2}\to\cdots
\end{equation}
の交差として表される。

\subsection{構造記述の限界}

\begin{proposition}[$F_0$ 指示限界原理]
数学的体系は、観測対象と観測者像の評価関係、文脈変換に沿う再構成、および時間点間の関係を記述できる。しかし、無点化された構造だけから現実の「ここ・いま・これ」を選出することは、構造記述とは別の指示作用を要する。
\end{proposition}

これは数学上の一般的不可能性定理ではない。$\CircleRef$ を追加データとして内部化することはできる。しかし、その内部点が現実の指示点であるという対応を確立するには、体系とその使用状況を結ぶ解釈が必要になる。本原理は、その解釈作用を構造そのものと混同しないためのメタ分析原理である。

\subsection{二つの視点の最終整理}

\begin{center}
\begin{tabular}{p{0.2\textwidth} p{0.34\textwidth} p{0.34\textwidth}}
\toprule
視点 & 問うこと & 形式的手がかり \\
\midrule
直交性 & 意味、指示、時間を型の異なる変化方向として分離できるか & 積、添字圏、依存型 \\
文脈的指示軸 & 指示位置と意味ファイバーをどう型分けするか & contextual category、C-system、再添字付け $f^*$ \\
四種類の時間 & 依存・計算・パス・観測の順序をどう区別するか & 文脈拡張、簡約、同一性型、点付き時間 \\
鏡像構成 & 対象と観測者像が評価を介してどう相互規定されるか & Chu双対、Isbell共役 \\
四重結像 & 四方向の自己同一性検査がどこで同時に閉じるか／崩れるか & QRB、四比較写像、四残余の共通台 \\
文脈再構成 & 観測文脈の変更後に像をどう延長するか & 左・右カン拡張 \\
現実的指示 & どの位置と時点が「ここ・いま・これ」か & 点付き構造、$\CircleRef$ \\
指示の先行性 & 観測意味の成立が何に依存するか & 指示先行公理、依存型族 $\mathcal M(r,t)$ \\
\bottomrule
\end{tabular}
\end{center}

直交性は分離の視点であり、contextual categoryはその指示方向を文脈と代入の基底軸として形式化する。鏡像構成は結合の視点であり、四重結像はその結合を四つの自己同一性検査へ分解する視点である。カン拡張は四つの像を文脈変換に沿って再構成し、それぞれが文脈を横断して保存されるかを問う。時間はそれらの観測事象を接続する。観測者と時間は $\CircleRef$ において交差するが、同一ではない。さらに $\CircleRef$ は、対称な鏡像関係の観測側を現実の指示原点として選び、観測された意味に対する指示の操作的先行性を成立させる。

\section{結論}

本稿は、意味と指示の関係を、直交性と鏡像構成という二つの相補的視点から再構成した。直交性は、意味／観測対象と指示／観測者像が異なる型の変化方向であることを示す。鏡像構成は、両方向が評価写像において結合され、転置によって相互の像を与えることを示す。したがって両者は別の現象ではなく、一つの評価構造を分解と結合の両面から捉えたものである。

指示方向については、Cartmellのcontextual categoryおよびVoevodskyのC-systemを導入し、文脈と代入の圏を基底、各文脈上の型と項を意味ファイバーとして配置した。これにより、指示の軸は比喩的な第二軸にとどまらず、$\Ctx$ 上の再添字付け $f^*$ として形式化される。依存和・依存積の随伴 $\Sigma_f\dashv f^*\dashv\Pi_f$ は局所的な文脈移送を、$\Lan_K\dashv K^*\dashv\Ran_K$ は関手圏上の大域的な文脈変更を与える。

時間構造については、依存時間、計算時間、パス時間、観測時間を区別した。依存時間は文脈形成、計算時間は項の実行、パス時間は同一性と変形、観測時間は現在点をもつ事象順序を表す。これらは追加構造によって関係づけられるが、自動的には同一視されない。特に文脈 $\Gamma$ と観測時点 $t$ は別の添字であり、意味構造は $\Obj:\Ctx^{\op}\times\Time\to\mathbf{Cat}$ として二重に相対化できる。

Chu空間は対象と観測文脈を評価写像の左右へ配置する静的模型を与え、Isbell双対は両側の表現を相互に共役させる。カン拡張は、観測文脈の変換後に成立する像を普遍的に延長または補完する。この三者を区別することで、鏡そのもの、鏡像の相互反射、文脈変更後の像の再構成を混同せずに記述できる。

QRBとの接続により、単一のChu鏡面は、同一性截断、二重否定閉包、随伴単位、自己参照比較という四つの結像様式へ分解された。四重結像は四つの像の一致ではなく、四比較写像を同じ対象上で重ねて検査する構造である。四比較がすべて可逆な状態は四重不動点を与え、それらが同じ局所点ですべて非可逆となる領域はQRBの外点候補領域を与える。真理保存性はこの全体の一断面であり、QRBの○を単独で導出するものではない。

この四重化によって、カン拡張は一つの像の延長ではなく、四つの結像様式が文脈変更と両立するかを検査する構造として精密化された。絶対カン拡張は、四自己関手を後続適用しても普遍性が保存される強い条件として位置づけられる。他方、四比較の共通障害領域を実際にQRBとして構成するには、自己関手とカン拡張の比較写像、余ファイバー、および台を追加しなければならない。

さらに、観測者と時間を分離した。観測者は区別の位置、時間は事象の接続構造である。両者は点付き観測 $\sigma:1\to\Obs\times\Time$ において「ここ・いま」として交差する。$F_0$ の $\bigcirc$ は、この議論では評価面 $\CircleEval$ と指示点 $\CircleRef$ に型分けされる。$\CircleEval$ は対象と観測者像を対称な鏡像関係へ置くのに対し、$\CircleRef$ は特定の観測位置を選び、その位置を観測意味の型付けに操作上先行させる。指示先行公理は意味の存在一般を観測へ還元せず、観測された意味の現実化に限定してこの依存関係を述べる。数学はこれらの構造と関係を記述できるが、現実的指示を無点化構造そのものと同一視しない。この非同一性を保持することが、意味と指示の構造論における $F_0$ の役割である。

今後の課題は、第一に $\CircleEval$ と $\CircleRef$ のあいだの型付き遷移を定義すること、第二にcontextual categoryまたはC-system上で意味ファイバー $\Obj(\Gamma)$ と指示先行公理を厳密に構成すること、第三に文脈相対化された四自己関手 $M_i^{\Gamma}$ と再添字付け $f^*$ の比較写像 $\beta_{i,f}$ を構成し、その余ファイバーの共通台を調べること、第四に観測文脈をプロ関手として構成しIsbell共役とカン拡張を統一すること、第五に時間添字をもつChu空間または高次圏的モデルを構成することである。これらが与えられたとき、本稿の構造的類比は、より具体的な数学的モデルへ進む。

\bibliographystyle{plain}
\bibliography{ref}

\end{document}